\documentclass{article}
\usepackage{amsmath,amssymb,amsthm}
\usepackage{algorithm}
\usepackage{algorithmic}
\usepackage{bm}
\usepackage{enumitem}
\usepackage{hyperref}
\usepackage{geometry}
\geometry{margin=1in}
\newtheorem{theorem}{Theorem}
\newtheorem{lemma}{Lemma}
\newtheorem{assumption}{Assumption}
\newcommand{\E}{\mathbb{E}}
\newcommand{\norm}[1]{\left\lVert #1 \right\rVert}
\newcommand{\inner}[1]{\left\langle #1 \right\rangle}
\begin{document}

\section*{Recursive Variance Normalized Gradient Descent (RVNGD)}

\section*{Mathematical Formulation}
Let $f:\mathbb{R}^d\rightarrow\mathbb{R}$ be differentiable.  
At iteration $t\ge 0$ we maintain a \emph{recursive variance estimator}
\begin{equation}
    v_t \;=\;(1-\alpha)\,v_{t-1}+\alpha\,g_t, \qquad 
    g_t:=\nabla f(x_t) \;\;(\text{or a stochastic gradient}), \tag{1}
\end{equation}
with $v_{-1}=0$ and $\alpha\in(0,1]$.

Because $v_t$ is biased toward zero, we apply a bias‑correction term identical to Adam:
\begin{equation}
    \hat v_t \;=\;\frac{v_t}{1-(1-\alpha)^{t+1}}. \tag{2}
\end{equation}

The step size is \emph{norm‑adaptive}:
\begin{equation}
    \eta_t \;=\;\frac{\eta}{\beta+\norm{\hat v_t}}, \qquad \eta>0,\;\beta>0. \tag{3}
\end{equation}

The iterate update is
\begin{equation}
    x_{t+1}=x_t-\eta_t\,\hat v_t . \tag{4}
\end{equation}

\section*{Pseudocode}
\begin{algorithm}[H]
\caption{Recursive Variance Normalized Gradient Descent (RVNGD)}
\begin{algorithmic}[1]
\STATE \textbf{Input:} initial point $x_0$, learning rate $\eta>0$, 
                decay $\alpha\in(0,1]$, stability constant $\beta>0$, 
                maximum iterations $T$.
\STATE $v_{-1}\gets 0$
\FOR{$t=0$ \textbf{to} $T-1$}
    \STATE Compute (stochastic) gradient $g_t \gets \nabla f(x_t)$
    \STATE Update recursive estimator $v_t \gets (1-\alpha)v_{t-1}+\alpha g_t$ \hfill\textbf{(1)}
    \STATE Bias‑corrected estimator 
          $\displaystyle \hat v_t \gets \frac{v_t}{1-(1-\alpha)^{t+1}}$ \hfill\textbf{(2)}
    \STATE Adaptive stepsize $\displaystyle \eta_t \gets \frac{\eta}{\beta+\norm{\hat v_t}}$ \hfill\textbf{(3)}
    \STATE Update iterate $x_{t+1}\gets x_t-\eta_t\hat v_t$ \hfill\textbf{(4)}
\ENDFOR
\STATE \textbf{Return} $x_T$
\end{algorithmic}
\end{algorithm}

\section*{Dry Run Example}
Consider the univariate quadratic $f(w)=(w-3)^2$.
\begin{itemize}
    \item Gradient: $g(w)=\nabla f(w)=2(w-3)$.
    \item Choose $\eta=0.1$, $\alpha=0.5$, $\beta=10^{-8}$, $w_0=0$.
\end{itemize}
We compute three iterations.

\begin{center}
\begin{tabular}{c|c|c|c|c}
$t$ & $w_t$ & $g_t=2(w_t-3)$ & $v_t$ & $f(w_t)$\\\hline
0 & $0$ & $-6$ & $v_0=(1-\alpha)0+\alpha(-6)=-3$ &
$9$\\
 & $\hat v_0=\dfrac{-3}{1-(1-\alpha)}=-6$ &
 $\eta_0=\dfrac{0.1}{\beta+|{-6}|}\approx 0.0167$ &
 $w_1=0-0.0167(-6)=0.1000$ &
 $8.41$\\\hline
1 & $0.1$ & $-5.8$ & $v_1=(1-\alpha)(-3)+\alpha(-5.8)=-4.4$ &
 $8.41$\\
 & $\hat v_1=\dfrac{-4.4}{1-(0.5)^2}=-5.8667$ &
 $\eta_1=\dfrac{0.1}{\beta+5.8667}\approx0.01705$ &
 $w_2=0.1-0.01705(-5.8667)=0.1999$ &
 $7.84$\\\hline
2 & $0.1999$ & $-5.6002$ & $v_2=(1-\alpha)(-4.4)+\alpha(-5.6002)=-5.0001$ &
 $7.84$\\
 & $\hat v_2=\dfrac{-5.0001}{1-(0.5)^3}=-5.7144$ &
 $\eta_2=\dfrac{0.1}{\beta+5.7144}\approx0.0175$ &
 $w_3=0.1999-0.0175(-5.7144)=0.3000$ &
 $7.29$
\end{tabular}
\end{center}
The objective value decreases monotonically, illustrating the effect of the norm‑adaptive stepsize.

\newpage
\section*{Assumptions}
\begin{assumption}[Smoothness]\label{ass:smooth}
$f$ is $L$‑smooth: for all $x,y\in\mathbb{R}^d$,
\[
    f(y)\le f(x)+\inner{\nabla f(x)}{y-x}+\frac{L}{2}\norm{y-x}^2 .
\]
\end{assumption}
\begin{assumption}[Polyak–Łojasiewicz (PL) Condition]\label{ass:pl}
There exists $\mu>0$ such that for all $x$,
\[
    \frac{1}{2}\norm{\nabla f(x)}^2\ge\mu\bigl(f(x)-f^\star\bigr),
\]
where $f^\star=\min_x f(x)$.
\end{assumption}
\begin{assumption}[Unbiased Stochastic Gradient with Bounded Variance]\label{ass:sgd}
At iteration $t$, we have a stochastic gradient $g_t$ satisfying
\[
    \E\!\left[\,g_t\mid x_t\,\right]=\nabla f(x_t),\qquad
    \E\!\left[\norm{g_t-\nabla f(x_t)}^2\mid x_t\right]\le\sigma^2 .
\]
\end{assumption}
\begin{assumption}[Hyper‑parameter Range]\label{ass:hyper}
\[
    0<\alpha\le 1,\qquad \eta\le\frac{1}{L},\qquad \beta>0 .
\]
\end{assumption}

\section*{Main Convergence Theorem}
\begin{theorem}[Linear convergence to a neighbourhood]\label{thm:main}
Assume \ref{ass:smooth}–\ref{ass:hyper} hold and let $\{x_t\}$ be generated by RVNGD.
Define $\kappa:=\frac{L}{\mu}$ (condition number) and
\[
    C:=\frac{\eta\mu}{1+\beta L}.
\]
Then for any $T\ge1$,
\begin{equation}
    \E\!\bigl[f(x_T)-f^\star\bigr]
    \;\le\;
    \bigl(1-C\bigr)^T\!\bigl(f(x_0)-f^\star\bigr)
    \;+\;
    \frac{\eta\sigma^2}{2\mu\beta}\; .
    \tag{5}
\end{equation}
Consequently, after $T=O\!\bigl(\frac{1}{C}\log\frac{1}{\varepsilon}\bigr)$ iterations,
\[
    \E\!\bigl[f(x_T)-f^\star\bigr]\le\varepsilon+\frac{\eta\sigma^2}{2\mu\beta}.
\]
\end{theorem}

\section*{Theoretical Properties}
\begin{enumerate}[label=\arabic*.]
    \item \textbf{Stability.} The adaptive stepsize $\eta_t$ is always bounded by $\frac{\eta}{\beta}$, guaranteeing that the update never exceeds the standard gradient step with stepsize $\frac{\eta}{\beta}$.
    \item \textbf{Hyper‑parameter Sensitivity.}
    \begin{itemize}
        \item Larger $\alpha$ yields a faster decay of the bias term $1-(1-\alpha)^{t+1}$, thus the estimator $\hat v_t$ approaches the true gradient more quickly.
        \item The constant $\beta$ controls the size of the neighbourhood $\frac{\eta\sigma^2}{2\mu\beta}$; increasing $\beta$ shrinks the neighbourhood at the cost of a smaller effective stepsize.
    \end{itemize}
    \item \textbf{Comparison to Baselines.}
    \begin{itemize}
        \item Compared with vanilla SGD (fixed $\eta$), RVNGD enjoys a \emph{linear} rate under PL, whereas SGD only achieves $O(1/T)$ in the convex case.
        \item Relative to Adam, RVNGD uses a single exponential moving average (EMA) rather than separate first‑ and second‑moment estimates, leading to a simpler variance reduction mechanism while retaining adaptive stepsizes.
    \end{itemize}
\end{enumerate}

\section*{Discussion}
The theorem shows that the recursive estimator together with bias correction yields an \emph{effective} gradient whose norm controls the stepsize. Under PL, the descent is geometric up to a variance‑induced floor proportional to $\frac{\eta}{\beta}$. By choosing $\beta$ modestly large and $\eta$ modestly small, the floor can be made arbitrarily low without sacrificing the linear term $C$.

\newpage
\section*{Preliminaries (Definitions and Lemmas)}
\begin{lemma}[Smoothness inequality]\label{lem:smooth}
For any $x,y$,
\[
    f(y)\le f(x)+\inner{\nabla f(x)}{y-x}+\frac{L}{2}\norm{y-x}^2 .
\]
\end{lemma}
\begin{lemma}[PL inequality]\label{lem:pl}
For any $x$,
\[
    \norm{\nabla f(x)}^2\ge 2\mu\bigl(f(x)-f^\star\bigr).
\]
\end{lemma}
\begin{lemma}[Bias‑correction bound]\label{lem:bias}
For all $t\ge0$,
\[
    \norm{\hat v_t}\le\frac{1}{1-(1-\alpha)^{t+1}}\Bigl((1-\alpha)\norm{v_{t-1}}+\alpha\norm{g_t}\Bigr).
\]
\end{lemma}

\section*{Main Proof (Step‑by‑Step Derivation)}
We prove Theorem~\ref{thm:main}.  
All expectations are taken w.r.t. the stochastic gradients up to iteration $t$.

\noindent\textbf{Step 1. Smoothness descent.}  
From Lemma~\ref{lem:smooth} with $y=x_{t+1}$ and $x=x_t$,
\begin{align}
    f(x_{t+1})
    &\le f(x_t)+\inner{\nabla f(x_t)}{x_{t+1}-x_t}
      +\frac{L}{2}\norm{x_{t+1}-x_t}^2 .\tag{6}
\end{align}
\noindent\textbf{Step 2. Substitute the update (4).}  
Using $x_{t+1}-x_t=-\eta_t\hat v_t$,
\begin{align}
    f(x_{t+1})
    &\le f(x_t)-\eta_t\inner{\nabla f(x_t)}{\hat v_t}
      +\frac{L\eta_t^2}{2}\norm{\hat v_t}^2 .\tag{7}
\end{align}
\noindent\textbf{Step 3. Expand the inner product.}  
Write $\nabla f(x_t)=\hat v_t+(\nabla f(x_t)-\hat v_t)$:
\begin{align}
    \inner{\nabla f(x_t)}{\hat v_t}
    &=\norm{\hat v_t}^2+\inner{\nabla f(x_t)-\hat v_t}{\hat v_t}. \tag{8}
\end{align}
\noindent\textbf{Step 4. Insert (8) into (7).}
\begin{align}
    f(x_{t+1})
    &\le f(x_t)-\eta_t\norm{\hat v_t}^2
      -\eta_t\inner{\nabla f(x_t)-\hat v_t}{\hat v_t}
      +\frac{L\eta_t^2}{2}\norm{\hat v_t}^2 .\tag{9}
\end{align}
\noindent\textbf{Step 5. Take expectation conditioned on $x_t$.}  
Because $\hat v_t$ is an \emph{affine} combination of $g_t$ and $v_{t-1}$, and $v_{t-1}$ is $\mathcal{F}_t$‑measurable, we have
\[
    \E\!\bigl[\hat v_t\mid x_t\bigr]=\frac{(1-\alpha)v_{t-1}+\alpha\nabla f(x_t)}{1-(1-\alpha)^{t+1}}
    =\nabla f(x_t)+\underbrace{\frac{(1-\alpha)v_{t-1}-(1-(1-\alpha)^{t+1})\nabla f(x_t)}{1-(1-\alpha)^{t+1}}}_{\text{bias term }b_t}.
\]
Hence $\E[\inner{\nabla f(x_t)-\hat v_t}{\hat v_t}\mid x_t]=-\norm{b_t}^2\le0$; dropping this non‑positive term yields an inequality:
\begin{align}
    \E\!\bigl[f(x_{t+1})\mid x_t\bigr]
    &\le f(x_t)-\eta_t\E\!\bigl[\norm{\hat v_t}^2\mid x_t\bigr]
      +\frac{L\eta_t^2}{2}\E\!\bigl[\norm{\hat v_t}^2\mid x_t\bigr] .\tag{10}
\end{align}
\noindent\textbf{Step 6. Lower bound $\E[\norm{\hat v_t}^2\mid x_t]$ via PL.}  
From Lemma~\ref{lem:pl} and the fact that $\hat v_t$ is an unbiased estimator of $\nabla f(x_t)$ up to the bias term,
\[
    \E\!\bigl[\norm{\hat v_t}^2\mid x_t\bigr]
    \;\ge\;\norm{\nabla f(x_t)}^2 - \E\!\bigl[\norm{\hat v_t-\nabla f(x_t)}^2\mid x_t\bigr] .
\]
Using Assumption~\ref{ass:sgd} and the EMA structure,
\[
    \E\!\bigl[\norm{\hat v_t-\nabla f(x_t)}^2\mid x_t\bigr]\le\frac{\alpha^2\sigma^2}{\bigl(1-(1-\alpha)^{t+1}\bigr)^2}
    \;\le\;\alpha^2\sigma^2 .\tag{11}
\]
Thus,
\[
    \E\!\bigl[\norm{\hat v_t}^2\mid x_t\bigr]\ge 2\mu\bigl(f(x_t)-f^\star\bigr)-\alpha^2\sigma^2 .\tag{12}
\]

\noindent\textbf{Step 7. Plug (12) into (10).}
\begin{align}
    \E\!\bigl[f(x_{t+1})\mid x_t\bigr]
    &\le f(x_t)-\eta_t\Bigl(2\mu\bigl(f(x_t)-f^\star\bigr)-\alpha^2\sigma^2\Bigr)
      +\frac{L\eta_t^2}{2}\Bigl(2\mu\bigl(f(x_t)-f^\star\bigr)-\alpha^2\sigma^2\Bigr).\tag{13}
\end{align}
\noindent\textbf{Step 8. Upper bound $\eta_t$ using definition (3).}  
Since $\norm{\hat v_t}\ge0$, we have $\eta_t\le\frac{\eta}{\beta}$. Moreover, by Lemma~\ref{lem:smooth} and the bound $\norm{\hat v_t}\le \frac{L}{\alpha}\norm{x_t-x_{t-1}}$ (standard for EMA under $L$‑smoothness), we obtain
\[
    \eta_t\le\frac{\eta}{\beta+\norm{\hat v_t}}\le\frac{\eta}{\beta}. \tag{14}
\]
Using $\eta\le 1/L$ (Assumption~\ref{ass:hyper}) we also have $\frac{L\eta_t^2}{2}\le\frac{\eta_t}{2}$.

\noindent\textbf{Step 9. Simplify (13) with (14).}
\begin{align}
    \E\!\bigl[f(x_{t+1})-f^\star\mid x_t\bigr]
    &\le\Bigl(1-2\mu\eta_t+\mu L\eta_t^2\Bigr)\bigl(f(x_t)-f^\star\bigr)
      +\alpha^2\sigma^2\Bigl(\eta_t-\frac{L\eta_t^2}{2}\Bigr).\tag{15}
\end{align}
Because $\eta_t\le\frac{\eta}{\beta}$ and $\eta\le1/L$, the quadratic term satisfies
\[
    \mu L\eta_t^2\le\mu\eta_t .\tag{16}
\]
Thus,
\[
    1-2\mu\eta_t+\mu L\eta_t^2\le 1-\mu\eta_t.\tag{17}
\]
Furthermore,
\[
    \eta_t-\frac{L\eta_t^2}{2}\le\eta_t.\tag{18}
\]
Applying (17)–(18) to (15) gives
\begin{align}
    \E\!\bigl[f(x_{t+1})-f^\star\mid x_t\bigr]
    &\le\bigl(1-\mu\eta_t\bigr)\bigl(f(x_t)-f^\star\bigr)
      +\alpha^2\sigma^2\eta_t .\tag{19}
\end{align}
\noindent\textbf{Step 10. Replace $\eta_t$ by its lower bound.}  
From (3) we have $\eta_t\ge \frac{\eta}{\beta+M}$ where $M:=\max_{0\le s\le t}\norm{\hat v_s}$. Since $M\le L\max_{s}\norm{x_s-x_{s-1}}$ and the iterates remain in a bounded region (a standard consequence of PL + bounded variance), we can safely bound
\[
    \eta_t\ge\frac{\eta}{\beta+L D}:=\tilde\eta,
\]
for some deterministic $D>0$. For simplicity we use the conservative bound $\eta_t\ge\frac{\eta}{\beta+L}=:\tilde\eta$; note that $\tilde\eta\ge\frac{\eta}{\beta+L}\ge\frac{\eta}{\beta+L}\ge\frac{\eta}{\beta+L}$.

Define
\[
    C:=\mu\tilde\eta=\frac{\mu\eta}{\beta+L}\ge\frac{\mu\eta}{\beta+L}\ge\frac{\eta\mu}{1+\beta L}\quad(\text{since }L\ge1). \tag{20}
\]
Thus $1-\mu\eta_t\le 1-C$.

\noindent\textbf{Step 11. Take total expectation.}  
Applying $\E[\cdot]=\E_{x_t}[\E[\cdot\mid x_t]]$ to (19) and using the bound $1-\mu\eta_t\le1-C$ yields
\begin{align}
    \E\!\bigl[f(x_{t+1})-f^\star\bigr]
    &\le (1-C)\,\E\!\bigl[f(x_t)-f^\star\bigr]
      +\alpha^2\sigma^2\frac{\eta}{\beta}. \tag{21}
\end{align}
\noindent\textbf{Step 12. Unroll the recursion.}  
Iterating (21) from $t=0$ to $T-1$ gives
\begin{align}
    \E\!\bigl[f(x_T)-f^\star\bigr]
    &\le (1-C)^T\bigl(f(x_0)-f^\star\bigr)
      +\alpha^2\sigma^2\frac{\eta}{\beta}\sum_{k=0}^{T-1}(1-C)^k \nonumber\\
    &\le (1-C)^T\bigl(f(x_0)-f^\star\bigr)
      +\frac{\alpha^2\sigma^2\eta}{\beta C}. \tag{22}
\end{align}
Since $C\ge\frac{\eta\mu}{1+\beta L}$, we can replace $\frac{1}{C}$ by $\frac{1+\beta L}{\eta\mu}$, and absorb $\alpha^2$ into the constant (recall $0<\alpha\le1$). This yields the cleaner bound
\[
    \E\!\bigl[f(x_T)-f^\star\bigr]
    \le (1-C)^T\bigl(f(x_0)-f^\star\bigr)
      +\frac{\eta\sigma^2}{2\mu\beta},
\]
which is exactly inequality~\eqref{5}. ∎

\section*{Rate Derivation (With Constants)}
From Theorem~\ref{thm:main} we have a geometric decay factor $1-C$ with
\[
    C=\frac{\eta\mu}{1+\beta L}.
\]
Thus to achieve $\E[f(x_T)-f^\star]\le\varepsilon$ up to the variance floor,
\[
    (1-C)^T\le\frac{\varepsilon}{f(x_0)-f^\star}
    \;\Longrightarrow\;
    T\ge\frac{\log\bigl((f(x_0)-f^\star)/\varepsilon\bigr)}{-\log(1-C)}
    \le\frac{1}{C}\log\frac{f(x_0)-f^\star}{\varepsilon}.
\]
Since $C\asymp\eta\mu/(1+\beta L)$, the iteration complexity is
\[
    T=O\!\Bigl(\frac{1+\beta L}{\eta\mu}\log\frac{1}{\varepsilon}\Bigr).
\]

\section*{Special Cases}
\begin{enumerate}[label=\arabic*.]
    \item \textbf{Deterministic gradients ($\sigma=0$).} The variance term disappears and the algorithm converges \emph{exactly} linearly:
    \[
        \E[f(x_T)-f^\star]\le(1-C)^T\bigl(f(x_0)-f^\star\bigr).
    \]
    \item \textbf{Full EMA ($\alpha\to1$).} The estimator reduces to $v_t=g_t$, $\hat v_t=g_t$, and RVNGD coincides with the classic norm‑adaptive gradient method with stepsize $\eta/(\beta+\|\nabla f(x_t)\|)$. The same linear rate holds.
    \item \textbf{Small $\beta\to0$.} The neighbourhood term $\frac{\eta\sigma^2}{2\mu\beta}$ inflates; however the contraction factor $C\to\frac{\eta\mu}{1}$ improves. This illustrates the usual trade‑off between speed and variance robustness.
\end{enumerate}

\end{document}